% =========================================================================
%  STRONA TYTUŁOWA
%
%  1. Uzupełnij dane w bloku poniżej – trafią one na stronę tytułową,
%     do oświadczenia i do metadanych PDF.
%  2. Część „Układ strony” nie wymaga edycji.
% =========================================================================

% ------------------------- DANE PRACY ------------------------------------
\autor{Imię Nazwisko}
\plecautora{M}              % M – formy męskie, K – żeńskie (w oświadczeniu)

\tytulPL{Temat pracy dyplomowej w~języku polskim}

\rodzajpracy{inżynierska}   % inżynierska / magisterska
\formastudiow{stacjonarne}  % stacjonarne / niestacjonarne
\stopienstudiow{I}          % I – studia inżynierskie, II – magisterskie
\kierunek{Inżynieria Środowiska}  % np. Energetyka, Geoinformatyka
\specjalnosc{brak}

\promotor{dr inż. Imię Nazwisko}
\recenzent{dr hab. inż. Imię Nazwisko}  % puste – pole „Recenzent” znika

\miejsce{Kraków}
\rok{2026}

\zgodaASAP{TAK}             % zgoda na udostępnienie pracy w ASAP: TAK / NIE
% -------------------------------------------------------------------------

\metadanepdf

% ------------------------- UKŁAD STRONY ----------------------------------
\newgeometry{top=1.7cm, bottom=1.3cm, left=2.2cm, right=2.2cm}
\thispagestyle{empty}
\begingroup
\setstretch{1}
\fontfamily{phv}\selectfont         % krój bezszeryfowy, jak na wzorze wydziału
\setlength{\parindent}{0pt}
\newcommand{\etykieta}[1]{{\fontsize{9.5}{11}\selectfont\color{black!55}#1}}

% Nagłówek: logo wydziału i sygnet PK
\logowydzialu\hfill\logopk\par
\vspace{5mm}
{\color{wisieniebieski}\rule{\linewidth}{0.6pt}}\par

\vspace*{\stretch{1.3}}
\begin{center}
  {\fontsize{22}{27}\selectfont\color{pkgranat}\dana{autor}\par}
  \vspace{9mm}
  {\fontsize{17}{22}\selectfont\bfseries\color{pkgranat}\dana{tytulPL}\par}
  \vspace{6mm}
  {\fontsize{11}{13}\selectfont praca \dana{rodzajpracy}\par}

  \vspace{13mm}
  {\fontsize{11}{17}\selectfont
    \textbf{studia \dana{formastudiow} \dana{stopienstudiow} stopnia}\\
    \etykieta{kierunek studiów:} \dana{kierunek}\\
    \etykieta{specjalność:} \dana{specjalnosc}\par}
\end{center}

\vspace*{\stretch{1}}
\begin{flushright}
  \etykieta{Promotor:}\\[2pt]
  {\fontsize{11.5}{14}\selectfont\bfseries\dana{promotor}}
  \jeslidana{recenzent}{\\[9mm]
    \etykieta{Recenzent:}\\[2pt]
    {\fontsize{11.5}{14}\selectfont\dana{recenzent}}}
\end{flushright}

\vspace*{\stretch{1.1}}
\begin{center}
  {\fontsize{11.5}{14}\selectfont\dana{miejsce}, \dana{rok}}
\end{center}
\vspace*{\stretch{0.5}}

% Stopka z danymi wydziału
{\color{wisieniebieski}\rule{\linewidth}{0.4pt}}\par
\vspace{2.5mm}
\begin{minipage}[b]{0.62\linewidth}
  \fontsize{7.5}{9.5}\selectfont
  {\bfseries\color{pkgranat}Wydział Inżynierii Środowiska i~Energetyki}\\
  ul. Warszawska 24, 31-155 Kraków\\
  tel. 12 628 28 01, s-0@pk.edu.pl
\end{minipage}\hfill
\begin{minipage}[b]{0.36\linewidth}
  \raggedleft\fontsize{11}{13}\selectfont\bfseries\color{wisieniebieski}%
  www.wisie.pk.edu.pl
\end{minipage}\par
\endgroup
\restoregeometry
\cleardoublepage
